% !TeX program = pdflatex
%==============================================================================
%  Lista Teórica 07 --- Classificação e Classificadores Gaussianos
%  O gabarito sai de "Lista teorica 07 - gabarito.tex".
%==============================================================================
\providecommand{\gabaritoopt}{}
\documentclass[11pt,a4paper]{article}
\usepackage[\gabaritoopt]{../../recursos/latex/estilo-lista}

\begin{document}
\cabecalholista{07}{Classificação e Classificadores Gaussianos}

%------------------------------------------------------------------------------
\begin{exercicio}[o classificador de Bayes, em uma dimensão]
Seja $Y\in\{0,1\}$ e, condicionalmente a $Y$,
\[
  X\mid Y=0 \sim N(-1,\,1), \qquad X\mid Y=1 \sim N(+1,\,1).
\]
\begin{enumerate}[label=(\alph*),leftmargin=1.1cm]
  \item Com $\Prob(Y=1)=\Prob(Y=0)=1/2$, mostre que o classificador de Bayes
        decide por $Y=1$ quando $x>0$.
  \item Calcule o \textbf{erro de Bayes} nesse caso, em função de $\Phi$ (a função
        de distribuição acumulada da normal padrão), e dê o valor numérico.
        \emph{Dado:} $\Phi(-1)=0{,}1587$.
  \item Agora suponha $\Prob(Y=1)=0{,}1$. Mostre que o corte passa a ser
        $x^\ast=\tfrac12\ln 9 \approx 1{,}0986$, e calcule o novo erro de Bayes.
        \emph{Dados:} $\Phi(-2{,}0986)=0{,}0179$ e $\Phi(0{,}0986)=0{,}5393$.
  \item Compare o resultado de (c) com o erro do classificador que responde
        sempre $Y=0$. O que isso antecipa sobre a Aula~08?
\end{enumerate}
\end{exercicio}

\begin{solucao}
\textbf{(a)} O classificador de Bayes escolhe a classe de maior probabilidade a
posteriori. Pelo Teorema de Bayes, com $f_0$ e $f_1$ as densidades de $X$ dado
$Y=0$ e dado $Y=1$ e $\pi_k=\Prob(Y=k)$,
\[
  \Prob(Y=1\mid X=x) = \frac{\pi_1 f_1(x)}{\pi_0 f_0(x)+\pi_1 f_1(x)},
\]
e o denominador, positivo, é o mesmo para as duas classes. Com duas classes, o
classificador de Bayes decide por $Y=1$ quando $\Prob(Y=1\mid X=x)>\Prob(Y=0\mid
X=x)$, isto é, quando $\pi_1 f_1(x) > \pi_0 f_0(x)$. Com $\pi_0=\pi_1$ as prioris
cancelam, e as constantes $1/\sqrt{2\pi}$ das duas densidades também; como a
exponencial é crescente,
\[
  \exp\Big(-\tfrac{(x-1)^2}{2}\Big) > \exp\Big(-\tfrac{(x+1)^2}{2}\Big)
  \iff -(x-1)^2 > -(x+1)^2
  \iff (x+1)^2 > (x-1)^2 .
\]
Expandindo, $x^2+2x+1 > x^2-2x+1 \iff 4x>0 \iff x>0$. O corte é o ponto médio
entre as duas médias, como a simetria já sugeria. Em $x=0$ as duas classes
empatam, $\Prob(Y=1\mid X=0)=\tfrac12$, e qualquer das duas respostas é ótima; como
$\Prob(X=0)=0$, a escolha nesse ponto não altera o risco.

\smallskip
\textbf{(b)} O classificador de Bayes é $g^\ast(x)=\1{x>0}$, e o seu risco é a
probabilidade de ele errar. Condicionando em $Y$ (lei da probabilidade total),
\[
  \risco^\ast = \tfrac12\,\Prob(X>0\mid Y=0) + \tfrac12\,\Prob(X\le 0\mid Y=1).
\]
Sob $Y=0$, $Z=X+1\sim N(0,1)$, então $\Prob(X>0\mid Y=0)=\Prob(Z>1)=1-\Phi(1)
=\Phi(-1)$. Sob $Y=1$, $Z=X-1\sim N(0,1)$, e $\Prob(X\le 0\mid Y=1)=\Prob(Z\le -1)
=\Phi(-1)$: as duas parcelas são iguais, como a simetria sugeria. Logo
\[
  \risco^\ast = \Phi(-1) = \mathbf{0{,}1587}.
\]
Quase 16\% de erro é o \emph{piso}: nenhum classificador, construído com quaisquer
dados, tem risco menor. É o análogo em classificação do $\sigma^2$ da Aula~01.

\smallskip
\textbf{(c)} Agora as prioris não cancelam. Decidimos por $Y=1$ quando
\[
  0{,}1\,\exp\Big(-\tfrac{(x-1)^2}{2}\Big) > 0{,}9\,\exp\Big(-\tfrac{(x+1)^2}{2}\Big).
\]
Tomando logaritmo,
\[
  \ln 0{,}1 - \tfrac{(x-1)^2}{2} > \ln 0{,}9 - \tfrac{(x+1)^2}{2}
  \iff \ln\tfrac{0{,}1}{0{,}9} > \tfrac{(x-1)^2-(x+1)^2}{2} = -2x,
\]
ou seja $-\ln 9 > -2x$, isto é $x > \tfrac12\ln 9 \approx 1{,}0986$.

O erro fica
\[
  \risco^\ast = 0{,}9\,\Prob(X>x^\ast\mid Y=0) + 0{,}1\,\Prob(X<x^\ast\mid Y=1)
   = 0{,}9\,\Phi(-2{,}0986) + 0{,}1\,\Phi(0{,}0986),
\]
usando $\Prob(X>x^\ast\mid Y=0)=1-\Phi(x^\ast+1)=\Phi(-x^\ast-1)$ e
$\Prob(X<x^\ast\mid Y=1)=\Phi(x^\ast-1)$. Numericamente,
\[
  \risco^\ast = 0{,}9\times 0{,}0179 + 0{,}1\times 0{,}5393
   = 0{,}01611 + 0{,}05393 = \mathbf{0{,}0700}.
\]

\smallskip
\textbf{(d)} O classificador que responde sempre $Y=0$ erra exatamente
$\Prob(Y=1)=\mathbf{0{,}10}$. O classificador de Bayes erra $0{,}0700$ --- melhor,
mas não muito: a diferença é de 3 pontos percentuais.

E vale olhar de onde vem o erro de Bayes: das duas parcelas, $0{,}0539$ vêm de
classificar como $0$ observações que eram $1$. Ou seja, o classificador ótimo
\textbf{erra mais da metade dos positivos} ($\Prob(X<x^\ast\mid Y=1)=0{,}5393$) ---
e ainda assim é ótimo. Mover o corte para a esquerda, de $x^\ast$ para $t<x^\ast$,
acertaria os positivos que caem em $(t,x^\ast]$, mas erraria os negativos que caem
no mesmo intervalo, e o risco mudaria de
\[
  \int_t^{x^\ast}\big[0{,}9\,f_0(x)-0{,}1\,f_1(x)\big]\,dx > 0,
\]
porque à esquerda de $x^\ast$ vale $0{,}9\,f_0(x)>0{,}1\,f_1(x)$ --- é a própria
definição do corte. Acertar mais positivos custaria mais falsos positivos, entre os
$90\%$ de negativos, do que os falsos negativos que se evitam.

É exatamente a situação em que a acurácia deixa de ser informativa, e é o assunto
da Aula~08. Quando a prevalência é desbalanceada, um número de acurácia próximo
do trivial pode esconder tanto o classificador ótimo quanto um inútil.
\end{solucao}

%------------------------------------------------------------------------------
\begin{exercicio}[por que a fronteira do LDA é linear]
No modelo gaussiano com covariância \textbf{comum}, as $K$ classes têm prioris
$\pi_k=\Prob(Y=k)$ e $\X\mid Y=k \sim N(\bm\mu_k,\,\bm\Sigma)$, com $\bm\Sigma$
simétrica e positiva definida; $f_k$ é a densidade de $\X$ dado $Y=k$.
\begin{enumerate}[label=(\alph*),leftmargin=1.1cm]
  \item Escreva $\ln\big(\pi_k f_k(\x)\big)$ e mostre que, ao comparar duas
        classes, os termos que dependem de $\x$ de forma \emph{quadrática}
        se cancelam.
  \item Conclua que o classificador de Bayes decide pela classe que maximiza a
        \emph{função discriminante linear}
        \[
          \delta_k(\x) = \x^\top\bm\Sigma^{-1}\bm\mu_k
            - \tfrac12\bm\mu_k^\top\bm\Sigma^{-1}\bm\mu_k + \ln\pi_k .
        \]
  \item Em que ponto exato do item (a) a hipótese de covariância comum é usada?
        O que sobra se cada classe tiver a sua, $\bm\Sigma_k$?
  \item No caso $p=1$, $\bm\Sigma=\sigma^2$, $\mu_0=-1$, $\mu_1=+1$ e
        $\pi_0=\pi_1$, recupere o corte do Exercício~1(a) a partir de
        $\delta_1(x)=\delta_0(x)$.
\end{enumerate}
\end{exercicio}

\begin{solucao}
\textbf{(a)} Primeiro, por que olhar para $\ln\big(\pi_k f_k(\x)\big)$. O
classificador de Bayes escolhe o $k$ que maximiza
$\Prob(Y=k\mid\X=\x)=\pi_k f_k(\x)/f(\x)$, em que $f(\x)=\sum_s\pi_s f_s(\x)$ não
depende de $k$; como o logaritmo é crescente, maximizar $\pi_k f_k(\x)$ ou o seu
logaritmo dá a mesma classe. A densidade gaussiana multivariada dá
\[
  \ln\big(\pi_k f_k(\x)\big)
  = \ln\pi_k - \tfrac{p}{2}\ln(2\pi) - \tfrac12\ln\det\bm\Sigma
    - \tfrac12(\x-\bm\mu_k)^\top\bm\Sigma^{-1}(\x-\bm\mu_k).
\]
Expandindo a forma quadrática, e usando que $\bm\Sigma^{-1}$ é simétrica --- então
os dois termos cruzados, $\x^\top\bm\Sigma^{-1}\bm\mu_k$ e
$\bm\mu_k^\top\bm\Sigma^{-1}\x$, são o mesmo número ---,
\[
  (\x-\bm\mu_k)^\top\bm\Sigma^{-1}(\x-\bm\mu_k)
  = \underbrace{\x^\top\bm\Sigma^{-1}\x}_{\text{não depende de }k}
    - 2\,\x^\top\bm\Sigma^{-1}\bm\mu_k
    + \bm\mu_k^\top\bm\Sigma^{-1}\bm\mu_k .
\]
Ao comparar duas classes $k$ e $\ell$, subtraímos as duas expressões. Os termos
$-\tfrac p2\ln(2\pi)$, $-\tfrac12\ln\det\bm\Sigma$ e
$-\tfrac12\x^\top\bm\Sigma^{-1}\x$ são \textbf{idênticos} nas duas e desaparecem
--- e é justamente o último deles que era quadrático em $\x$.

\smallskip
\textbf{(b)} Descartando os termos comuns a todas as classes --- que somam a mesma
quantidade a cada uma e, por isso, não mudam qual delas é a maior ---, sobra
\[
  \delta_k(\x) = \x^\top\bm\Sigma^{-1}\bm\mu_k
    - \tfrac12\bm\mu_k^\top\bm\Sigma^{-1}\bm\mu_k + \ln\pi_k ,
\]
em que o primeiro termo vem de
$-\tfrac12\cdot\big(-2\,\x^\top\bm\Sigma^{-1}\bm\mu_k\big)$. Ela é \textbf{afim em
$\x$}: um produto interno com o vetor fixo $\bm\Sigma^{-1}\bm\mu_k$, mais uma
constante. As classes $k$ e $\ell$ empatam no conjunto
$\{\x: \delta_k(\x)=\delta_\ell(\x)\}$, isto é,
\[
  \x^\top\bm\Sigma^{-1}(\bm\mu_k-\bm\mu_\ell)
  = \tfrac12\big(\bm\mu_k^\top\bm\Sigma^{-1}\bm\mu_k
    - \bm\mu_\ell^\top\bm\Sigma^{-1}\bm\mu_\ell\big) + \ln\frac{\pi_\ell}{\pi_k},
\]
que é um \textbf{hiperplano} sempre que $\bm\mu_k\ne\bm\mu_\ell$: como
$\bm\Sigma^{-1}$ é invertível, o vetor normal $\bm\Sigma^{-1}(\bm\mu_k-\bm\mu_\ell)$
só se anula quando as médias coincidem. Com duas classes, esse hiperplano é a
fronteira; com mais, a fronteira entre $k$ e $\ell$ é o pedaço dele em que as duas
empatam e superam as demais.

\smallskip
\textbf{(c)} A hipótese é usada exatamente ao afirmar que
$\x^\top\bm\Sigma^{-1}\x$ e $\ln\det\bm\Sigma$ \emph{não dependem de $k$}. Com
covariâncias distintas teríamos $\x^\top\bm\Sigma_k^{-1}\x$ e
$\ln\det\bm\Sigma_k$, que dependem de $k$ e \textbf{não se cancelam}: se
$\bm\Sigma_k\ne\bm\Sigma_\ell$, então $\bm\Sigma_k^{-1}\ne\bm\Sigma_\ell^{-1}$, e a
diferença $\x^\top\big(\bm\Sigma_k^{-1}-\bm\Sigma_\ell^{-1}\big)\x$ não é
identicamente nula.

Sobra então
\[
  \delta_k(\x) = -\tfrac12\x^\top\bm\Sigma_k^{-1}\x
    + \x^\top\bm\Sigma_k^{-1}\bm\mu_k
    - \tfrac12\bm\mu_k^\top\bm\Sigma_k^{-1}\bm\mu_k
    - \tfrac12\ln\det\bm\Sigma_k + \ln\pi_k ,
\]
que é \emph{quadrática} em $\x$ --- é o QDA. A fronteira entre duas classes passa a
ser uma quádrica; em $\R^2$, uma cônica (elipse, parábola, hipérbole ou um caso
degenerado delas, como um par de retas). Em uma frase: \textbf{dentro do modelo
gaussiano, o que torna a fronteira linear não é a normalidade, é a covariância
comum} --- com covariâncias diferentes, as mesmas normais dão fronteira quadrática.

\smallskip
\textbf{(d)} Com $p=1$, $\bm\Sigma^{-1}=1/\sigma^2$ e $\pi_0=\pi_1$ (então
$\ln\pi_k$ cancela):
\[
  \delta_1(x)=\frac{x\cdot 1}{\sigma^2}-\frac{1}{2\sigma^2},
  \qquad
  \delta_0(x)=\frac{x\cdot(-1)}{\sigma^2}-\frac{1}{2\sigma^2}.
\]
Igualando: $x/\sigma^2 = -x/\sigma^2$, ou seja $2x/\sigma^2=0$, isto é
$\mathbf{x=0}$. É o corte do Exercício~1(a), e note que, com prioris iguais, ele
\emph{não depende de $\sigma^2$} --- a variância muda o quanto se erra, não onde fica
a fronteira. Com prioris diferentes o $\ln\pi_k$ não cancela, e
$\delta_1(x)=\delta_0(x)$ dá $x=\tfrac{\sigma^2}{2}\ln\frac{\pi_0}{\pi_1}$: aí a
variância também desloca o corte. Com $\sigma^2=1$, $\pi_0=0{,}9$ e $\pi_1=0{,}1$, é
o $x^\ast=\tfrac12\ln 9$ do Exercício~1(c).
\end{solucao}

%------------------------------------------------------------------------------
\begin{exercicio}[quantos parâmetros cada um estima]
Considere $K=2$ classes em $\R^p$.
\begin{enumerate}[label=(\alph*),leftmargin=1.1cm]
  \item Conte os parâmetros livres do LDA e do QDA. Considere: as médias das
        classes, a(s) matriz(es) de covariância --- lembrando que uma matriz
        simétrica $p\times p$ tem $p(p+1)/2$ entradas livres --- e as prioris.
  \item Preencha a tabela para $p=2$, $p=10$ e $p=50$.
  \item Mostre que a diferença QDA menos LDA vale $p(p+1)/2$ e comente a ordem de
        grandeza.
  \item A tabela abaixo traz a acurácia do LDA e do QDA em $p=2$ e $p=10$, em
        função do tamanho $n$ do treino (média sobre 60 amostras, avaliada em 20
        mil observações novas). Nas duas dimensões, as classes são gaussianas com
        covariâncias diferentes: o QDA é o modelo correto. O traço marca onde o
        QDA não pôde ser ajustado. Nas duas o QDA passa à frente entre $n=30$ e
        $n=50$ --- mas o que acontece \emph{antes} disso é muito diferente. Use a
        contagem de (b) para explicar, inclusive o traço, e diga por que a regra
        ``prefira LDA com poucos dados'' está incompleta.
\end{enumerate}

\begin{center}\small
\renewcommand{\arraystretch}{1.2}
\begin{tabular}{@{}llrrrrrr@{}}
\toprule
        &     & $n=20$ & $30$ & $50$ & $100$ & $300$ & $1000$ \\
\midrule
\multirow{2}{*}{$p=2$}  & LDA & $0{,}8602$ & $0{,}8681$ & $0{,}8724$ & $0{,}8775$ & $0{,}8785$ & $0{,}8793$ \\
                        & QDA & $0{,}8508$ & $0{,}8669$ & $0{,}8737$ & $0{,}8806$ & $0{,}8824$ & $0{,}8834$ \\
\midrule
\multirow{2}{*}{$p=10$} & LDA & $0{,}7053$ & $0{,}7422$ & $0{,}7773$ & $0{,}8013$ & $0{,}8175$ & $0{,}8241$ \\
                        & QDA & --- & $0{,}7030$ & $0{,}7896$ & $0{,}8440$ & $0{,}8792$ & $0{,}8905$ \\
\bottomrule
\end{tabular}
\end{center}
\end{exercicio}

\begin{solucao}
\textbf{(a)} Com $K=2$:
\begin{itemize}[leftmargin=1.4cm]
  \item \textbf{médias}: dois vetores de $\R^p$, logo $2p$ parâmetros, nos dois
        métodos;
  \item \textbf{covariância}: o LDA estima \emph{uma} matriz comum,
        $p(p+1)/2$; o QDA estima \emph{duas}, $2\cdot p(p+1)/2 = p(p+1)$;
  \item \textbf{prioris}: $\pi_0$ e $\pi_1$ somam 1, então há $K-1=1$ parâmetro
        livre.
\end{itemize}
\[
  \#\text{LDA} = 2p + \tfrac{p(p+1)}{2} + 1,
  \qquad
  \#\text{QDA} = 2p + p(p+1) + 1 .
\]

\smallskip
\textbf{(b)}
\begin{center}\small
\renewcommand{\arraystretch}{1.3}
\begin{tabular}{@{}lrrr@{}}
\toprule
$p$ & $2$ & $10$ & $50$ \\
\midrule
LDA & $8$  & $76$  & $1\,376$ \\
QDA & $11$ & $131$ & $2\,651$ \\
\midrule
diferença & $3$ & $55$ & $1\,275$ \\
\bottomrule
\end{tabular}
\end{center}

\smallskip
\textbf{(c)} A diferença é
$\big(2p+p(p+1)+1\big)-\big(2p+\tfrac{p(p+1)}{2}+1\big)=\tfrac{p(p+1)}{2}$: a
segunda matriz de covariância, exatamente. Ela cresce como $p^2/2$, ou seja
\textbf{quadraticamente}. Em $p=2$ são 3 parâmetros a mais --- nada. Em $p=50$ são
$1\,275$ a mais, quase o tamanho do modelo inteiro do LDA ($1\,376$): o QDA fica
com quase o dobro dos parâmetros.

\smallskip
\textbf{(d)} Na grade medida, a virada cai no mesmo intervalo nas duas dimensões
--- entre $n=30$ e $n=50$ --- e justamente por isso ela é a parte \emph{menos}
informativa da tabela. O que muda é o \textbf{tamanho} da vantagem de cada lado:

\begin{center}\small
\renewcommand{\arraystretch}{1.2}
\begin{tabular}{@{}lcc@{}}
\toprule
                                    & $p=2$ & $p=10$ \\
\midrule
LDA $-$ QDA em $n=30$ (antes da virada)   & $+0{,}0012$ & $\mathbf{+0{,}0392}$ \\
QDA $-$ LDA em $n=1000$ (depois)          & $+0{,}0041$ & $\mathbf{+0{,}0664}$ \\
custo da matriz extra, $p(p+1)/2$         & $3$ & $55$ \\
\bottomrule
\end{tabular}
\end{center}

Em $p=10$ a escolha importa muito, nos dois sentidos. O traço em $n=20$ é a
contagem de (b) levada ao extremo: são $10$ observações por classe para estimar,
em cada uma, uma matriz com $55$ parâmetros. Pior: a covariância amostral de $10$
pontos em $\R^{10}$ tem posto no máximo $9$, porque os dez vetores, centrados na
média da classe, somam zero e ficam num subespaço de dimensão no máximo $9$. Ela
não é invertível, e o QDA, que precisa de $\widehat{\bm\Sigma}_k^{-1}$ e de
$\ln\det\widehat{\bm\Sigma}_k$, não está definido --- só uma versão regularizada,
que puxa $\widehat{\bm\Sigma}_k$ na direção de uma matriz invertível, poderia ser
ajustada. Com $n=30$, quinze observações por classe, a covariância já é invertível
e o QDA ajusta, mas fica $3{,}9$ pontos atrás do LDA. Com $n=1000$ ele fica a menos
de um ponto do classificador de Bayes, que acerta $0{,}8964$ nessa população, e
$6{,}6$ pontos à frente do LDA.

Em $p=2$ os dois métodos ficam a menos de um ponto percentual um do outro em toda
a tabela. A matriz extra custa só $3$ parâmetros, então o QDA tem pouco a
\emph{perder}. E, nesta população, ele também tem pouco a \emph{ganhar}: o
classificador de Bayes acerta $0{,}8841$, e o LDA ajustado com um milhão de
observações, $0{,}8805$ --- praticamente a melhor reta possível, que acerta
$0{,}8806$. Nem com infinitos dados o QDA, que converge para o classificador de
Bayes, abriria mais que $0{,}4$ ponto sobre o LDA. Não é que as covariâncias sejam
iguais: as correlações são parecidas ($-0{,}47$ e $-0{,}45$), mas a variância de
$X_2$ é $3{,}85$ numa classe e $1{,}83$ na outra. É que, onde as classes se
misturam, essa diferença pouco curva a fronteira de Bayes. E isso não é uma
propriedade de $p=2$: numa população em $\R^2$ com correlações de sinais opostos,
o QDA ganha do LDA com folga desde $n=20$ --- é o que mede a \texttt{Aula prática
07}.

Portanto a regra ``prefira LDA com poucos dados'' está incompleta em dois
sentidos. Primeiro, \textbf{o que decide não é $n$ sozinho, e sim $n$ contra o que
o QDA estima a mais}: $p(p+1)/2$ parâmetros, e mais de $p$ observações por classe
para que cada covariância amostral seja sequer invertível --- um custo que só pesa
quando $p$ é grande. Segundo, a regra fala só do que o QDA tem a perder:
\textbf{o que ele tem a ganhar é a diferença entre o erro do LDA com dados
infinitos e o erro de Bayes}, ou seja, o quanto a fronteira ótima se afasta de um
hiperplano onde há dados. Isso nenhuma contagem de parâmetros informa --- é
assunto para a validação cruzada.

É a mesma lógica do balanço viés--variância da Aula~01: o QDA é o modelo mais
flexível, e modelos flexíveis só compensam quando há dados para sustentá-los e
estrutura para capturar.
\end{solucao}

\end{document}
